\begin{abstract}
Inverse Synthetic Aperture Radar (ISAR) imaging is advantageous to surveillance and reconnaissance efforts as it enables imaging in any weather, day or night. However, the resolution of the image is dependent on angular diversity and is susceptible to errors from the higher-order motion parameters of a non-cooperative target. For example, if the pulse repetition frequency (PRF) is insufficient to process the incoming Doppler bandwidth of the received echoes, the image can have aliased scatterers. Additionally, target scatterers rarely fit to a defined grid typically used for ISAR imaging, which creates matching basis errors. Though existing methods in literature have demonstrated good reconstruction quality, many have inherent limitations. Learning-based methods require training which can be difficult for ISAR imaging as publicly available real data is scarce. Other methods, like sparse Bayesian learning (SBL) and atomic norm minimization (ANM), are computationally burdensome which can be disadvantageous for imaging platforms with hardware and software limitations or in scenarios where real-time imaging is paramount. Our proposed solution, a modified Orthogonal Matching Pursuit (OMP) algorithm, is able to use the non-uniform motion of the target to aid in disambiguating target scatterers, and our algorithm efficiently finds off-grid target scatterers through Newton’s method. Our algorithm identifies aliased targets in an unambiguous region and determines the originating ambiguity by assessing the correlation of each ambiguous atom with the residual. Our modified OMP processes the latent image in a singular ambiguity as opposed to expanding the sensing matrix to incorporate other ambiguities. This reduces the sensing dictionary dimensionality. From the estimated scattering position, our algorithm then applies Newton’s method in a local neighborhood to refine the off-grid position. Preliminary results demonstrate the algorithm’s effectiveness relative to other modified OMP algorithms, outperforming each in positional accuracy and computational burden.
\end{abstract}


% Include a list of keywords after the abstract 
\keywords{ISAR, doppler ambiguity, non-uniform motion, off-grid}
